\section{Introduction}

This section sets the context for the research by outlining the underlying motivation, the proposed methodological approach, and the boundaries of the study. It guides the reader from the broader societal challenge of urban climate adaptation to the specific technical implementation of tree detection using satellite imagery.

\subsection{Motivation and Project Overview}

As urban areas expand and global temperatures rise due to climate change, being able to understand how trees can help cool cities has become increasingly important \cite{eu2023tree}. Furthermore, trees reduce the risk of flooding in cities and promote health and wellbeing \cite{woodlandtrust2023}. Our project aims at taking the first step in the direction of better tree coverage in urban areas by creating a model that is able to predict precise tree locations using satellite imagery of densely populated areas. \\
\\
This research builds upon the foundation laid by previous projects conducted in collaboration with the partner company and environmental simulation platform \textit{infrared.city}. A recent flood risk analysis developed for the platform highlighted that surface impermeability is a dominant factor in urban flooding \cite{marco2026floodrisk}. However, the authors noted that the lack of detailed, highly localized environmental data limits the precision of urban risk analyses \cite{marco2026floodrisk}. Concurrently, a second project successfully mapped general urban green spaces using Random Forest models on Sentinel-2 data \cite{gotschim2026greenspace}. While effective for broad classification, that project grouped all vegetation into a single category and reported that initial attempts to use U-Net architectures for detailed spatial analysis failed to yield good results \cite{gotschim2026greenspace}. The authors explicitly recommended future work to revisit U-Net architectures and to distinguish vegetation types more effectively \cite{gotschim2026greenspace}. \\
\\
Recently, there has been a rapidly growing body of research on urban tree detection using deep learning models. One of the most recent papers is by Gong et al. \cite{gong2024treedetector}, who use satellite data not only to detect trees but also to estimate their size and density. Our project tries to build on both this literature and the identified shortcomings of the previous \textit{infrared.city} projects. By integrating multi-season imagery and temporal attention algorithms, we aim to overcome the previous limitations of the U-Net model and achieve a more granular, object-specific classification than general green spaces. 

\subsection{Proposed Approach and Data}

To achieve reliable urban tree detection, our approach combines multi-spectral satellite imagery with a deep learning architecture optimized for image segmentation. \\
\\
We utilize publicly available Sentinel-2 satellite images, providing a 10-meter spatial resolution across four color channels: red, green, blue, and near-infrared. These images are retrieved via the API of the Copernicus Programme and are provided by the European Space Agency \cite{copernicus2025}. A key technical focus of our data strategy is the inclusion of imagery from different seasons to account for the significant impact of seasonality on tree appearance. To provide ground truth for training and testing, we integrate the publicly available ``Baumkataster'' dataset, which contains exact locations for every registered tree in the city of Vienna \cite{baumkataster2025}. \\
\\
Regarding the model, we employ a U-Net architecture, a specialized Convolutional Neural Network (CNN) consisting of symmetric contracting and expanding paths \cite{aramendia2024}. While a U-Net implementation in a previous project did not yield promising results, this paper aims to build upon that foundation as the architecture is considered highly accurate for detecting objects in satellite data \cite{pesek2024convolution, garnot2021panoptic}. Our model extracts features from seasonal multi-channel imagery through its contracting path and utilizes an expanding path for upsampling to predict precise locations. \\
\\
To improve upon previous attempts, our implementation incorporates skip connections and a temporal attention algorithm to prioritize more meaningful images within the time series \cite{garnot2021panoptic}. The model performance is optimized using Binary Cross-Entropy loss, comparing our predictions against the split training and test sets derived from the Vienna ``Baumkataster''.

\subsection{Scope and Limitations}

The model in this paper is trained only on satellite images of urban areas and is therefore only supposed to be used for tree detection within city environments. Additionally, the model focuses on tree detection in the temperate climate zone, as this is the vegetation it is trained to observe. Further analyses on the impact of trees on, for example, heat regulation in cities and flooding prevention are out of the scope of this project and could be the topic of future research. \\
\\
Since this paper uses Sentinel-2 satellite images with 10-meter spatial resolution, issues could arise, as the images might not have a high enough resolution to detect trees reliably. Future work could therefore also focus on improving image quality by using better available satellite data.

\subsection{Source Code Availability}
The source code, datasets, and scripts used to produce the results in this project are open source and can be found in the following GitHub repository: \\
\url{https://github.com/nikmaster/wu_dslab_infrared-city}
